
\subsection*{4.7.}

On
revient à la situation de 3.7, où on avait le diagramme commutatif :
\[
\begin{tikzcd}[column sep=large,row sep=large]
Y' \arrow[r,"j"] \arrow[d,"g"'] & X' \arrow[d,"f"]\\
Y \arrow[r,"i"] & X,
\end{tikzcd}
\]
où \(i\) est une immersion régulière, \(X'\) le \(X\)-schéma obtenu en faisant éclater \(Y\), et \(Y'=X'\times_XY\). On suppose \(X\) noethérien avec un faisceau ample, et \(Y\) régulier. On a vu en 3.7 qu'on a alors la suite exacte :
\[
0\longrightarrow K^{\bullet}(Y)\xrightarrow{u}
K^{\bullet}(Y')\times K^{\bullet}(X)\xrightarrow{v}
K^{\bullet}(X')\longrightarrow0.
\]
Nous allons montrer que cette suite exacte permet de décrire toute la situation formée par les \(K^{\bullet}\) de \(X,Y,X',Y'\) à partir de la donnée des \(\lambda\)-anneaux \(K^{\bullet}(X)\), \(K^{\bullet}(Y)\), des applications \(i^*\) et \(i_*\), et de l'élément \(N\) de \(K^{\bullet}(Y)\). (du moins à torsion près, jusqu'à nouvel ordre).

Soient donc \(K,K'\) deux \(\lambda\)-anneaux, \(i_*:K\to K'\) un homomorphisme additif, \(i^*:K'\to K\) un \(\lambda\)-homomorphisme, \(d\) un entier, et \(N\) un élément de \(K\) tel que \(\lambda^i(N)=0\) pour \(i>d\) et que \(\lambda^d(N)\) soit inversible. On suppose de plus vraies les propriétés suivantes :
\begin{enumerate}[label=\roman*)]
\item \(i_*\) est \(K'\)-linéaire (« formule de projection »);
\item \(i^*\circ i_*\) est l'homomorphisme « multiplication par \(\lambda_{-1}(N)\) »;
\item pour tout \(n\ge1\), et tout \(y\in K\), \(i_*\bigl(\lambda^n(N,y)\bigr)=\lambda^n(i_*(y))\) (relation (4.3.1)
\end{enumerate}

On utilise d'abord VI.I.1 pour construire une \(K\)-\(\lambda\)-algèbre \(K_1\) qui joue le rôle de \(K^{\bullet}(Y')\). Soit \(L\) une indéterminée, et \(K[L]\) la \(K\)-\(\lambda\)-algèbre engendrée par le générateur \(L\) soumis aux relations \({}^i(L)=0\) pour \(i\ge2\) (V 4.9). Soit :
\[
  P(L)=\lambda^d(N-L)=\sum_{i=0}^{d}(-1)^i\lambda^{d-i}(N)L^i.
\]
